\chapter{Literature Review}
\label{chap:literature}

This chapter reviews the work on which HepatoTwin builds and the gaps it
addresses. It follows the structure of the system: the clinical problem and
the evidence for dietary treatment (Section~\ref{sec:lit_masld}); how liver fat
is measured without biopsy (Section~\ref{sec:lit_measurement}); machine-learning
models that estimate steatosis from routine data, and the limits on how well
they can do so (Section~\ref{sec:lit_ml}); how diet is measured and
personalised (Section~\ref{sec:lit_diet}); digital twins in health and
nutrition (Section~\ref{sec:lit_twins}); computational gastronomy and food
recommendation (Section~\ref{sec:lit_food}); and how a twin can be evaluated
when no longitudinal data are available (Section~\ref{sec:lit_insilico}).
Section~\ref{sec:lit_gap} summarises the gap.

% ======================================================================
\section{MASLD and Its Dietary Management}
\label{sec:lit_masld}

\subsection{Definition and burden}

Metabolic dysfunction-associated steatotic liver disease (MASLD) is defined as
hepatic steatosis together with at least one of five cardiometabolic risk
factors, in the absence of another cause \cite{rinella2023nomenclature}. The
2023 multi-society nomenclature replaced non-alcoholic fatty liver disease
(NAFLD), removed the reliance on exclusion and stigmatising terms, and
introduced MetALD for patients with MASLD who also drink moderate amounts of
alcohol. A meta-analysis of studies using imaging to diagnose NAFLD estimated
a global adult prevalence of 30.05\%, rising from 25.26\% in 1990--2006 to
38.00\% in 2016--2019 \cite{younossi2023global}; a second meta-analysis
reached a similar overall estimate \cite{riazi2022prevalence}. Most patients
have simple steatosis, but a minority progress to steatohepatitis, fibrosis
and cirrhosis, and the disease is closely associated with type~2 diabetes and
cardiovascular risk \cite{easl2024masld,rinella2023aasld}.

\subsection{Evidence for lifestyle and dietary treatment}

Both the European (EASL-EASD-EASO) and American (AASLD) guidelines name
lifestyle change as the first-line treatment for most patients
\cite{easl2024masld,rinella2023aasld}. The strongest evidence concerns weight
loss. In a prospective study of 293 patients with paired liver biopsies,
weight loss through lifestyle change produced dose-dependent histological
improvement: losses of 5\% or more improved steatosis and inflammation, and
losses of 10\% or more resolved steatohepatitis in 90\% of patients and led to
fibrosis regression in 45\% \cite{vilargomez2015weight}. A meta-analysis of 22
randomised trials with 2{,}588 participants found that weight-loss
interventions lowered ALT by 9.81~U/L and improved liver fat and histological
activity, but found no significant change in fibrosis
\cite{koutoukidis2019weightloss}. The thresholds of 5\%, 7--10\% and 10\%
weight loss are now written into the European guideline \cite{easl2024masld}.
Liver fat also responds quickly: in people with type~2 diabetes on an
800~kcal per day diet, hepatic triglyceride fell from 12.8\% to 2.9\% over eight
weeks, and had already fallen by 30\% after the first week
\cite{lim2011reversal}.

The evidence for specific dietary components is more mixed, and this
variation matters for any system that turns evidence into advice. A 2026
meta-analysis of randomised trials, which graded the certainty of each finding
with GRADE, found that a Mediterranean-style diet reduced MRI-PDFF by 1.37
percentage points and improved liver stiffness, while omega-3 supplementation
and low-carbohydrate high-fat diets did not change ALT \cite{stern2026dietary}.
Dietary interventions including calorie restriction improved liver enzymes,
steatosis and stiffness in another meta-analysis
\cite{haigh2022mediterranean}. Yet a meta-analysis of ten trials found that
the Mediterranean diet reduced AST and GGT but not ALT
\cite{sangouni2022mediterranean}, and it was no better than a low-fat diet in
the short term \cite{xiong2024mediterranean}.

Feeding trials show that the type of surplus energy matters. When 38
overweight adults were overfed by 1{,}000~kcal per day for three weeks, liver
fat rose by 55\% with saturated fat, 33\% with simple sugars and 15\% with
unsaturated fat \cite{luukkonen2018saturated}. A meta-analysis of 13 trials
with 260 participants found that fructose raised liver fat and ALT when it
added surplus energy, but had no effect when it replaced other carbohydrate in
the same amount of energy \cite{chiu2014fructose}. Exercise also reduces liver
fat: in a meta-analysis of trials measuring liver fat by MRI, exercise was
associated with a relative reduction of 30\% or more, the threshold that
predicts histological response, even without clinically significant weight loss
\cite{stine2023exercise}.

Two lessons follow. First, the evidence is graded: biopsy-based dose--response
data for weight loss are stronger than meta-analyses of liver enzymes, and null
findings sit alongside positive ones. A system that recommends food should
therefore carry the strength of the evidence with each recommendation, rather
than treating all claims as equal. Second, the evidence is expressed in terms
of energy, weight change, nutrients and dietary patterns, never single recipes;
any mapping from evidence to a recipe is an inference that must be made
explicit.

% ======================================================================
\section{Non-Invasive Assessment of Liver Fat}
\label{sec:lit_measurement}

Liver biopsy is the reference standard for grading steatosis and fibrosis but
is invasive and impractical for screening. Magnetic resonance imaging
proton-density fat fraction (MRI-PDFF) and proton spectroscopy
($^1$H-MRS) measure liver fat accurately but are costly. Vibration-controlled
transient elastography (FibroScan) is cheaper and widely used: its controlled
attenuation parameter (CAP), reported from 100 to 400~dB/m, grades steatosis,
with cut-offs of 248, 268 and 280~dB/m for grades S1, S2 and S3 derived from an
individual patient data meta-analysis \cite{karlas2017cap}.

CAP is, however, an imperfect measure of liver fat. A review of the technique
reports that agreement between two operators examining the same patient has a
concordance correlation of 0.82, falling to 0.65 in patients with BMI of
30~kg/m$^2$ or more and to 0.44 when CAP is below 240~dB/m
\cite{ferraioli2021cap}. The same review notes that CAP accuracy falls when the
interquartile range (IQR) of the ten acquisitions exceeds 40~dB/m, which is
why that limit is used as a quality criterion. In adults with NAFLD, CAP
correlated with MRI-PDFF at $r = 0.66$, falling to $r = 0.37$ for changes over
time \cite{ahn2018cap}. Part of every CAP reading is therefore measurement
variation, and the variation is largest in exactly the people who are hardest
to examine. This limits how closely any model can predict CAP
(Section~\ref{sec:lit_ml_bounds}).

Simple clinical scores offer a cheaper alternative. The Fatty Liver Index
(FLI) combines BMI, waist circumference, triglycerides and GGT
\cite{bedogni2006fli}; the Hepatic Steatosis Index (HSI) combines the ALT/AST
ratio, BMI, sex and diabetes \cite{lee2010hsi}; and the NAFLD liver fat score
combines metabolic syndrome, type~2 diabetes, fasting insulin, AST and the
AST/ALT ratio, reaching an AUC of 0.86--0.87 against spectroscopy, with genetic
information adding less than 1\% \cite{kotronen2009prediction}. These scores are
the natural baselines for any new steatosis model.

% ======================================================================
\section{Machine-Learning Prediction of Steatosis and Liver Fat}
\label{sec:lit_ml}

\subsection{Classification of steatosis in NHANES}

The US National Health and Nutrition Examination Survey (NHANES) added
FibroScan to its 2017--March 2020 cycle, and several studies have since
trained classifiers to detect CAP-defined steatosis from routine data
(Table~\ref{tab:lit_nhanes}). Gradient-boosted trees, usually XGBoost
\cite{chen2016xgboost}, perform best or near best in most of them, and the
strongest predictors are consistently body-size measures (waist circumference
and BMI), followed by liver enzymes, glucose and insulin markers, and
triglycerides. Model explanations, typically SHAP values
\cite{lundberg2017shap}, are widely reported.

Du et al.\ used 6{,}814 participants and found logistic regression, XGBoost
and LightGBM comparable (AUC 0.815--0.842), so that the added complexity of
boosting bought little \cite{du2026masld}. Liu et al.\ compared models on an
automated machine-learning platform in 4{,}177 participants; XGBoost reached an
AUC of 0.859, and the authors imputed missing values by multiple imputation and
reported calibration curves and decision-curve analysis
\cite{liu2024automl}. Zhang et al.\ reported an AUC of 0.874 on 2{,}460
complete cases \cite{zhang2025masld}. Boullion et al.\ identified biomarkers of
steatosis and stiffness; propensity matching on age and sex required complete
data and reduced the steatosis sample to 2{,}714, so their models were trained
on the larger unmatched data, where random forests reached 79.6\% accuracy
\cite{boullion2025biomarker}. Zhu et al.\ trained a three-grade model on
215{,}145 non-obese Chinese adults, reaching a macro-average AUC of 0.941
internally and 0.874 in 726 non-Hispanic Asian NHANES participants with
complete data \cite{zhu2026nonobese}; this is one of the few studies with
external validation, and it shows the usual fall in performance when a model
moves to a new population.

\begin{table}[htbp]
    \centering
    \small
    \caption{Machine-learning classifiers of CAP-defined steatosis using NHANES.}
    \label{tab:lit_nhanes}
    \begin{tabular}{p{3.0cm}p{3.6cm}p{2.6cm}p{2.4cm}p{2.6cm}}
        \hline
        Study & Population & Outcome & Best model & Performance \\
        \hline
        Du et al.\ (2026) \cite{du2026masld} & 6{,}814 adults, 2017--2020 & MASLD, CAP $\geq$ 280 & LR, XGBoost, LightGBM & AUC 0.815--0.842 \\
        Liu et al.\ (2024) \cite{liu2024automl} & 4{,}177, fasting subsample; imputed & CAP $\geq$ 302 & XGBoost (AutoML) & AUC 0.859; calibration curve \\
        Zhang et al.\ (2025) \cite{zhang2025masld} & 2{,}460 complete cases & CAP-defined MASLD & XGBoost & AUC 0.874 \\
        Boullion et al.\ (2025) \cite{boullion2025biomarker} & Unmatched sample (matched subset 2{,}714) & CAP $\geq$ 238 & Random forest & Accuracy 79.6\% \\
        Zhu et al.\ (2026) \cite{zhu2026nonobese} & 215{,}145 non-obese Chinese adults; 726 NHANES Asians & CAP grades (S0, S1, S2--S3) & XGBoost & Macro AUC 0.941 internal; 0.874 external \\
        \hline
    \end{tabular}
\end{table}

Three limitations recur. Most studies use only the 2017--2020 cycle, and most
discard participants with any missing value. Calibration is reported only
occasionally \cite{liu2024automl}, although a risk model whose output will be
shown to patients must produce probabilities that mean what they say;
isotonic regression is a standard way to obtain them from tree ensembles
\cite{zadrozny2002calibration}. And all are single snapshots, used to classify
a participant once rather than to track or support them over time.

\subsection{Prediction of liver fat as a continuous value}

Fewer studies predict the amount of liver fat, and their results are more
sobering (Table~\ref{tab:lit_regression}). In 28{,}796 UK Biobank adults, the
Fatty Liver Index explained 28\% of the variance in MRI-PDFF (33--37\% across
cross-validation folds); an extended
model with twelve variables (FLI+) explained 54\% on average across four
cross-validation folds (0.520--0.557) \cite{asaturyan2022fli}. A UK Biobank
preprint reported $R^2$ of 0.37 from anthropometric and blood measures in
40{,}533 participants with MRI-PDFF and 0.34 from plasma metabolites in 9{,}723,
against 0.63--0.64 from DXA body-composition images analysed by a
convolutional network \cite{insitro2024ukb}. Haas et al., also in UK Biobank, found that BMI
correlated with MRI liver fat at $r = 0.42$ and that adding 24 clinical
factors and biomarkers raised the correlation only to $r = 0.58$
($R^2 \approx 0.34$), concluding that liver fat cannot be reliably estimated
from clinical factors \cite{haas2021liverfat}. In Finnish adults measured by
spectroscopy, liver enzymes alone explained 15--25\% of the variance; an
equation including a specialist marker (phosphorylated IGFBP-1) reached an
adjusted $R^2$ of 0.49 in validation and an AUROC of 0.84 for NAFLD, against
0.72 for FLI and 0.62 for HSI \cite{petaja2016igfbp}. Higher values come from
small, selected samples: in 272 patients of a Chinese obesity clinic, a
four-variable equation predicted CAP with $R^2 = 0.76$ in derivation and
$0.65$ in validation \cite{liu2022surrogate}, and a four-variable index
correlated with MRI-PDFF at $\rho = 0.74$ in 149 overweight or obese adults,
validated only by bootstrap resampling \cite{chen2023shsi}. No study reporting
$R^2$ for CAP in NHANES was found.

\begin{table}[htbp]
    \centering
    \small
    \caption{Prediction of liver fat as a continuous value.}
    \label{tab:lit_regression}
    \begin{tabular}{p{3.3cm}p{3.1cm}p{2.0cm}p{3.8cm}p{2.3cm}}
        \hline
        Study & Population & Reference & Inputs & $R^2$ \\
        \hline
        UK Biobank preprint (2024) \cite{insitro2024ukb} & 40{,}533; 9{,}723; about 46{,}000 UK adults & MRI-PDFF & Body and blood measures; metabolites; DXA images & 0.37; 0.34; 0.63--0.64 \\
        Asaturyan et al.\ (2022) \cite{asaturyan2022fli} & 28{,}796 UK adults & MRI-PDFF & FLI; FLI with 12 variables (FLI+) & 0.28--0.36; 0.54 (4-fold mean) \\
        Haas et al.\ (2021) \cite{haas2021liverfat} & UK adults, 1{,}214 test & MRI & BMI; BMI with 24 clinical factors & $r$ 0.42; $r$ 0.58 \\
        Pet{\"a}j{\"a} et al.\ (2016) \cite{petaja2016igfbp} & 378 Finnish adults & $^1$H-MRS & Enzymes only; equation with phosphorylated IGFBP-1 & 0.15--0.25; 0.49 (adjusted) \\
        Liu et al.\ (2022) \cite{liu2022surrogate} & 272, obesity clinic & CAP & BMI, neck circumference, TG, C-peptide & 0.76 / 0.65 \\
        Chen et al.\ (2023) \cite{chen2023shsi} & 149, overweight or obese & MRI-PDFF & ALT, AST, insulin, 1-h glucose & $\rho = 0.74$ \\
        Kotronen et al.\ (2009) \cite{kotronen2009prediction} & 470 Finnish adults & $^1$H-MRS & Metabolic and liver markers & Not reported \\
        \hline
    \end{tabular}
\end{table}

\subsection{Why prediction from routine data is bounded}
\label{sec:lit_ml_bounds}

Two findings explain why continuous liver fat is hard to predict from routine
measurements. First, liver fat is strongly heritable: in a prospective study
of monozygotic and dizygotic twins using MRI-PDFF, the heritability of
steatosis was 0.52 after adjustment for age, sex and ethnicity
\cite{loomba2015heritability}, and in UK Biobank measured genetic variants
together explained up to 30\% of the variance in liver fat
\cite{haas2021liverfat}. Variants such as those in PNPLA3 change liver fat at
the same body size, and routine clinical data contain none of this
information. Second, the reference measurement itself is noisy
(Section~\ref{sec:lit_measurement}): a model can explain no more of the
variation in a single CAP reading than the reading shares with the patient's
true liver state. Together, these findings imply that a model built on routine
data should reach a moderate $R^2$ for liver fat, between about 0.3 and 0.55
in the large studies above, and good discrimination of steatosis, which is the
pattern seen in Tables~\ref{tab:lit_nhanes} and~\ref{tab:lit_regression}. For
a digital twin, the implication is that a calibrated probability of steatosis
is a more defensible state variable than a single predicted CAP value.

% ======================================================================
\section{Dietary Assessment and Personalised Nutrition}
\label{sec:lit_diet}

Personalising dietary advice requires knowing what a person usually eats.
NHANES collects one or two 24-hour dietary recalls per participant. A single
recall describes the average intake of a group well, but an individual's
usual intake poorly, because intake varies from day to day. Statistical
methods developed at the US National Cancer Institute estimate usual intake
from two or more recalls by separating within-person from between-person
variation \cite{dodd2006usual}. For foods that are not eaten every day, the
NCI method uses a two-part model that separates the probability of eating a
food on a given day from the amount eaten on days it is eaten, allows the two
parts to be correlated, and can add covariates such as age, sex or a food
frequency questionnaire; in simulation it estimated the distribution of usual
intake substantially better than earlier methods \cite{tooze2006usual}.
Self-reported intake is also biased: in the OPEN study, comparison with
recovery biomarkers showed substantial under-reporting, particularly of energy
\cite{subar2003open}. Diet quality can be summarised with the Healthy Eating
Index-2020, a 0--100 score of 13 components aligned with the 2020--2025 Dietary
Guidelines for Americans; its components are unchanged from HEI-2015 for ages
2 and over, and a separate index was introduced for toddlers
\cite{shamswhite2023hei}. Computing it requires food-group data from the USDA
Food Patterns Equivalents Database.

Personalisation is motivated by individual variation in response. Zeevi et
al.\ showed in 800 people that post-meal glucose responses to identical meals
differ widely and can be predicted from personal data, and that diets
personalised on these predictions lowered post-meal glucose
\cite{zeevi2015personalized}. For a liver twin, these findings suggest two
requirements: dietary inputs should represent usual rather than single-day
intake, and the system should become more personal as more of the patient's
own data arrive.

% ======================================================================
\section{Digital Twins in Health and Nutrition}
\label{sec:lit_twins}

\subsection{Definitions and requirements}

The US National Academies define a digital twin as a set of virtual
information constructs that mimics the structure, context and behaviour of a
physical system, is dynamically updated with data from it, has a predictive
capability and informs decisions; the bidirectional flow between the virtual
and physical system is central, and credibility rests on verification,
validation and uncertainty quantification (VVUQ) \cite{nasem2024digitaltwins}.
A scoping review of health twins frames them as a physical entity, its virtual
representation and the connection between them, and found that the term is
used for systems ranging from static models to fully coupled ones
\cite{katsoulakis2024digital}. A perspective on medical twins describes them as
most advanced in cardiology and oncology, and names data integration,
validation and privacy as the main obstacles \cite{laubenbacher2024medicine}.

Two architectural proposals are relevant to a modular twin. Nitschke et al.\
propose a general patient twin built on a bipartite knowledge graph that links
patient attributes to existing predictive models; fusion models based on
ensemble learning combine the models' outputs, and a propagation scheme passes
information through the graph each time new patient data arrive. They describe
the design as modular, evolving, informed, interpretable and explainable
\cite{nitschke2026design}. Nye structures patient data as a knowledge graph
built on medical ontologies such as SNOMED CT and proposes closed-form
continuous-time (``liquid'') neural networks to model the patient's evolving
state \cite{nye2023knowledge}. Both are conceptual designs without empirical
evaluation, but they share a principle that HepatoTwin adopts: separate
modules, each with a defined role, connected by an explicit patient state.

\subsection{Metabolic and nutrition twins}

In nutrition, Gkouskou et al.\ extended the idea of an ``actual digital
twin'', in which a patient is matched to many similar patients whose outcomes
are known, into a ``virtual digital twin'': a personal model combining
genetic, longitudinal metabolomic, immune, microbial and behavioural data,
collected through connected devices and managed by machine learning, to
simulate dietary interventions for that person \cite{gkouskou2020virtual}.
Metabolic twins have also been outlined as tools that track a person's
metabolic flexibility over time \cite{su2025metabolic}. A systematic review of
twins in precision health identified 12 studies and found improvements in most
of the reported outcomes, while noting that the evidence is still small and
heterogeneous \cite{shen2024effectiveness}.

The most developed clinical example is the Twin Health platform. Its method
paper describes a machine-learning model (CatBoost) that predicts each
person's post-meal glucose response from continuous glucose monitoring, food
logs, activity and sleep, and a mobile application that adjusts dietary
recommendations from those predictions and keeps refining them as data arrive
\cite{twinhealth2024frontiers}. In a one-year randomised trial in 319 people
with type~2 diabetes, nutrition guided by this twin improved liver-fat scores
and MRI-PDFF compared with standard care \cite{joshi2023digitaltwin}. A second
example, iDT-diet, modelled how diet quality affects biomarkers over time in
402 participants measured at baseline and after one, three and six years. It
combined a random forest with Choquet-integral LASSO feature selection, which
had lower error than XGBoost, LightGBM and CatBoost for several targets, with
generated natural-language descriptions and three-dimensional avatars of the
predicted body state \cite{nobel2026idtdiet}. In MASLD specifically, the
European ARTEMIS project is building virtual twins from retrospective data on
several thousand patients \cite{artemis2024}.

These systems either recommend individual foods from predicted glucose
responses or forecast biomarker trajectories. The Twin Health platform depends
on continuous glucose monitoring and proprietary data, and iDT-diet forecasts
biomarkers without recommending food. None evaluates complete, existing
recipes against evidence-graded dietary priorities, and none combines health
with cuisine and practicality when making a recommendation.

\subsection{From prediction to ``what if''}

A twin that compares interventions must answer counterfactual questions,
which predictive accuracy alone does not license. Freeman et al.\ name two
traps. The \emph{fidelity trap} is the belief that a model accurate enough to
reproduce patient trajectories can therefore say what would happen under a
different treatment; prediction and counterfactual estimation are distinct
tasks with different justifications. The \emph{feedback trap} arises when a
twin learns from data that its own recommendations helped to produce, which
breaks the assumptions behind its causal estimates. They argue that health
twins should be causally valid, modular and evolving systems, with modularity
isolating the parts that need interventional justification
\cite{freeman2026traps}. Laudy formalises a twin as a simulator of each
person's potential outcomes under treatment and control, and proposes a
hierarchy of fidelity (marginal, joint and structural) with a five-level
validation regime that turns the claim ``the simulator produces correct
counterfactuals'' into tests against observable data, while stating that the
dependence between a person's two potential outcomes remains untestable
\cite{laudy2026counterfactual}.

For HepatoTwin, these arguments have a direct consequence. A cross-sectional
model trained on observational data, such as a steatosis classifier built on
NHANES, describes the current state well but cannot by itself predict the
effect of changing a diet. Intervention effects must instead come from
interventional evidence, held in a separate module with their strength and
uncertainty made explicit, and the twin's updating must be evaluated in a
setting where the true effect is known (Section~\ref{sec:lit_insilico}).

% ======================================================================
\section{Computational Gastronomy and Food Recommendation}
\label{sec:lit_food}

\subsection{Recipe and food resources}

Computational gastronomy treats food and cooking as structured data, studying
recipes through their flavour, nutrition, health associations and
sustainability \cite{bagler2024gastronomy}. RecipeDB contains 118{,}171 recipes
from six continents, 26 geocultural regions and 74 countries, cooked with 268
processes from over 20{,}000 ingredients, with cuisine labels, utensils and
nutritional profiles, and ingredients linked to their flavour molecules
\cite{batra2020recipedb}. FlavorDB catalogues 25{,}595 flavour molecules, 2{,}254
of which are linked to 936 natural ingredients in 34 categories
\cite{garg2018flavordb}. Literature mining of 5{,}769 abstracts produced 8{,}957
associations between 152 culinary spices and 848 diseases; for 150 of the 152
spices the beneficial associations outnumbered the adverse ones, and the
authors used this balance to make disease-specific spice recommendations
\cite{rakhi2018spices}. These resources make it possible to reason about
recipes at scale, but they are built for exploration and analysis rather than
for deciding which recipes suit a particular patient.

\subsection{Food recommender systems}

Food recommendation faces challenges beyond those of other domains: it must
balance taste, health, culture and practicality. In a survey of the field,
Trattner and Elsweiler showed that re-weighting standard recommender outputs by
WHO or UK Food Standards Agency health scores can trade accuracy for
healthiness, but the healthiest lists came with very poor accuracy, and they
named health-aware recommendation an open problem \cite{trattner2017food}.

Knowledge graphs are a common foundation. FoodKG links over a million
recipes, their ingredients and USDA nutrient data in 67 million triples, and
supports queries such as which recipe can be made from the ingredients at hand
while avoiding allergens \cite{haussmann2019foodkg}. Chen et al.\ framed
personalised recommendation as constrained question answering over FoodKG,
treating the user's query, dietary preferences and health guidelines as
constraints, and improved on non-personalised baselines by 59.7 percentage
points on average across metrics \cite{chen2021constrained}. Graph neural
networks have been used to move users gradually towards healthier recipes
while respecting their preferences \cite{li2022healthguided}. In all of these
systems, ``healthy'' is defined by generic nutritional guidelines rather than
by a patient's clinical state, and the evidence behind each recommendation is
not graded.

\subsection{Recipe modification and adaptation}

A related line of work modifies recipes by substituting ingredients. Shirai et
al.\ combined explicit knowledge from a food knowledge graph with word
embeddings into a substitutability heuristic (DIISH), filtered the candidates
by nutritional and food-category constraints, and evaluated it against
substitutions taken from cooking guides and recipe reviews
\cite{shirai2021substitutions}. Loesch et al.\ learned knowledge-graph
embeddings (RDF2Vec) over a food graph and evaluated them against 1{,}847 known
substitution pairs, retrieving about 40\% of the true substitutes in the top
ten candidates \cite{loesch2022substitution}. Rita et al.\ fine-tuned large
language models to choose substitutions that raise phytochemical content,
improving Hit@1 from 34.5\% to 38.0\% on the GISMo benchmark and producing
1{,}639 enriched recipes, while cautioning that the health claims rest on
preclinical evidence \cite{rita2024phytochemical}. A survey of the area
concludes that a good substitute must match flavour, culinary function and
health suitability together, and recommends knowledge graphs and
neuro-symbolic reasoning for doing so \cite{kim2025substitution}. Cultural
adaptation goes further: translating recipes between Chinese and English
cuisines requires changing ingredients, units, tools and techniques, and
GPT-4 still fell short of human experts when adapting English recipes into
Chinese \cite{cao2024cultural}.

Every modification creates a new recipe whose taste, feasibility and nutrition
must be predicted, and none of these studies tested whether people liked or
cooked the result. Evaluating existing recipes, the approach taken in this
thesis, avoids that uncertainty: every recommended recipe has already been
written, cooked and nutritionally described.

\subsection{Menu planning, diversity and additional criteria}

Menu planning composes meals that satisfy nutritional and cost targets; Sahin
et al.\ released EvoMeal, an open-source tool comparing four multi-objective
evolutionary algorithms, and found that two newer algorithms (AGE-MOEA and
SMS-EMOA) outperformed NSGA-II and NSGA-III \cite{sahin2024menu}.
Recommendation lists also benefit from diversity: maximal marginal relevance
re-ranks items by a linear combination of relevance and dissimilarity to the
items already chosen, reducing redundancy \cite{carbonell1998mmr}. Two further
criteria are relevant to MASLD guidance. The degree of food processing can be
estimated with FPro, a machine-learning score that predicts processing from
nutrient profiles; it classified over 73\% of the US food supply as
ultra-processed, was associated in NHANES with metabolic syndrome and
diabetes, and showed that swapping foods for less-processed alternatives
could reduce these associations \cite{menichetti2023processing}. Glycaemic
load is available from the international glycaemic index tables
\cite{atkinson2021gi}.

% ======================================================================
\section{Evaluating Digital Twins Without Longitudinal Data}
\label{sec:lit_insilico}

A twin's defining features, updating with new data and projecting the future,
can only be fully validated against repeated measurements of the same people.
When such data are unavailable, simulation offers a recognised alternative.
In-silico clinical trials test interventions on computer-simulated patients
before, or alongside, real trials \cite{pappalardo2019insilico}. The best-known
example is the UVA/Padova type~1 diabetes simulator, a population of virtual
patients built from physiological models, which the US Food and Drug
Administration accepted in 2008 as a substitute for animal trials in the
preclinical testing of closed-loop insulin control
\cite{kovatchev2009insilico,dallaman2014simulator}. Its authors stress that
such testing is a prerequisite to clinical trials, not a substitute for them
\cite{kovatchev2009insilico}. Synthea generates synthetic patient records for
software testing at scale \cite{walonoski2018synthea}.

Two designs are particularly relevant. In plasmode simulation, real data
provide the covariates and the outcome is simulated from a known model, so the
data keep their real structure while the true effect is known and a method's
ability to recover it can be measured \cite{franklin2014plasmode}. In data
assimilation, observing system simulation experiments generate a ``true''
state with one model, give a forecasting system only noisy observations of it,
and measure how well the system recovers the truth; the field calls these twin
experiments \cite{kalnay2003assimilation}.

Simulated follow-up is credible only if its ingredients come from evidence.
For diet-driven weight change, Hall et al.\ developed and validated a dynamic
model of how a sustained change in energy intake translates into weight over
time: a lasting reduction of about 10~kcal per day eventually lowers weight by
about 0.45~kg, with half of the change reached in about one year and 95\% in
about three years. The common static rule of 3{,}500~kcal per pound
overestimates long-term weight loss by about 100\%, and the model predicts
average rather than individual responses \cite{hall2011energy}. Large
lifestyle trials report the weight trajectories that real participants
achieve. In Look AHEAD, intensive lifestyle intervention produced a mean
weight loss of 8.6\% at one year against 0.6\% with usual care; 68.0\% and
37.7\% of participants lost at least 5\% and 10\%, and at eight years the mean
losses were 4.7\% and 2.1\% \cite{lookahead2014eightyear}. The Diabetes
Prevention Program reported a similar pattern of loss and partial regain
\cite{dpp2002}, and a Look AHEAD analysis showed that losses of 5--10\%
increased the odds of clinically meaningful improvements in glycaemia, blood
pressure and lipids \cite{wing2011modest}. The response of the liver itself is
given by the weight-loss studies of Section~\ref{sec:lit_masld}
\cite{koutoukidis2019weightloss,vilargomez2015weight,lim2011reversal}.
Realistic observations also require the normal biological variation of each
laboratory test \cite{eflm_bv} and the known under-reporting of food intake
\cite{subar2003open}.

These sources make it possible to evaluate a twin's updating on virtual
patients whose follow-up is grounded in published human data. Such an
evaluation verifies that the twin behaves correctly when the truth is known;
it does not establish its accuracy for real patients, which still requires
longitudinal data.

% ======================================================================
\section{Summary and Research Gap}
\label{sec:lit_gap}

Table~\ref{tab:lit_gap} summarises what each strand of the literature offers
and what it leaves open. Clinical research provides graded evidence on how
weight and diet change the liver, but expressed in energy, nutrients and
patterns rather than recipes. Machine-learning studies estimate steatosis well
from routine data, but as one-off classifications, rarely calibrated, and with
a ceiling on continuous prediction set by genetics and measurement. Digital
twins in nutrition personalise food choices or forecast biomarkers, but do not
evaluate existing recipes against evidence-graded priorities, and the causal
literature warns that predictive accuracy alone does not justify their
intervention advice. Food recommender and substitution systems handle
preferences and generic healthiness, but not a patient's clinical state or the
strength of evidence. And no liver twin has been evaluated for how it updates
over time.

\begin{table}[htbp]
    \centering
    \small
    \caption{Strands of related work, what they leave open, and how HepatoTwin addresses it.}
    \label{tab:lit_gap}
    \begin{tabular}{p{3.0cm}p{4.0cm}p{3.8cm}p{4.0cm}}
        \hline
        Strand & What it provides & What it leaves open & HepatoTwin \\
        \hline
        Clinical evidence (Section~\ref{sec:lit_masld}) & Graded effects of weight loss, diet patterns and nutrients & Not expressed for recipes; evidence strength varies & Evidence engine with explicit strength labels, mapped to recipe criteria \\
        Steatosis models (Section~\ref{sec:lit_ml}) & Good discrimination from routine data & Snapshots; rarely calibrated; bounded $R^2$ & Calibrated probabilities as the twin's state, within a defined scope \\
        Dietary assessment (Section~\ref{sec:lit_diet}) & Recall data; usual-intake methods & Single-day recalls used directly & Recall-adjusted priorities; usual intake updated from food logs \\
        Digital twins (Section~\ref{sec:lit_twins}) & Personal state, glucose-response prediction, trajectories & No recipe-level, evidence-graded recommendation; fidelity and feedback traps & Modular twin linking state, evidence and recipes; effects from interventional evidence \\
        Food recommendation (Section~\ref{sec:lit_food}) & Recipe data, preferences, generic healthiness, substitution & No clinical state; ungraded evidence; untested modified recipes & Evidence-gated ranking of existing recipes by health, cuisine and feasibility \\
        In-silico evaluation (Section~\ref{sec:lit_insilico}) & Virtual patients, plasmode and twin experiments & Not yet applied to liver twins & In-silico evaluation of the twin's updating on virtual NHANES patients \\
        \hline
    \end{tabular}
\end{table}

This thesis addresses the intersection of these strands: a patient-specific
liver twin whose calibrated state, combined with graded intervention evidence
and the patient's own diet, drives the transparent evaluation and ranking of
existing recipes, and whose updating is evaluated in silico before real
longitudinal data exist.
